\noindent Protein language model (PLM) zero-shot variant-effect scores are often treated as general predictors despite task-specific evidence and incompatible score transforms. We audit raw ESM-2 log-likelihood ratios across six axes, with uncertainty where the evidence permits it and applicability-gated controls. Deep mutational scanning (DMS) discrimination rose to a protein-mean area under the receiver operating characteristic curve (AUROC) of 0.725 at 650M parameters, with no detected gain at 3B (paired difference $-0.00015$, 69 proteins, 95\% confidence interval, CI, $[-0.01340,0.01683]$); ProtBERT reached 0.711. Recomputed scores for 4,677 ClinVar variants give a mean AUROC of \ClinicalMacroAUC{} across 20 genes (gene-resampling CI $[\ClinicalMacroLow,\ClinicalMacroHigh]$). On the same scores, five-fold Platt and isotonic calibration yield Brier scores \PlattBrier{} and \IsoBrier{}, compared with \GenePrevBrier{} for a fitting-fold gene-prevalence baseline. Identity-grouped refitting supports these within-panel differences but does not establish transport to new genes or populations. The likelihood readout was a weak stability proxy. Saved antibody--antigen summaries give Spearman approximately 0.08 ($n=627$), but their interval provenance is not fully reconciled and does not support a significance claim. No Benjamini--Hochberg-significant association was detected for polyreactivity (0/14 datasets) or aggregation (0/7). These results support a bounded task-specific audit, not a universal reliability rule or a clinically deployable probability model.
